\section{Introduction}\label{sec:introduction}
This report documents a simulator of rocket flight in the vertical plane, with room left for a
later 3D extension, and the physics behind it. The simulator integrates a rigid rocket from the
launch rail to landing, so that design choices --- motor, mass, aerodynamics, recovery --- can be
compared by their effect on the trajectory. The report derives the equations the simulator solves
and records the models and assumptions behind each term.

Around the simulation core (\texttt{sim/}) sit the inputs and the tools that use it:
\begin{itemize}
    \item \textbf{Configurations and data.} Rocket, launch site, weather, motor and scenario
    descriptions (\texttt{configs/}), with motor curves, aerodynamic tables and atmospheric
    soundings under \texttt{data/}. The simulator itself is offline and deterministic; a sounding
    can be fetched separately with \texttt{tools/fetch\_sounding.py}.
    \item \textbf{Simulation and comparisons} (\texttt{run.py}): flies a scenario and compares
    variants of it.
    \item \sloppy\textbf{Applications} (\texttt{apps/}): motor selection (\texttt{run\_motor\_sweep.py}),
    air-brake sizing (\texttt{run\_airbrake.py}), selection of CFD conditions
    (\texttt{run\_cfd\_conditions.py}) and the mass--aerodynamics trade-off
    (\texttt{run\_trade.py}).
\end{itemize}

The report is organised as follows. \autoref{sec:rocket2d} sets up the frames, forces and equations
of motion; \autoref{sec:models} gives the interchangeable models behind each term;
\autoref{sec:analysis} works the equations through by hand and derives the stability results.
\autoref{sec:assumptions} lists the simplifications in force and \autoref{sec:backlog} what is left
undone. The appendix holds the notation, the model inventory and the sign conventions.
